\section{Methods}

\subsection{Study I system development}

Viscereality maps respiratory movement onto expansion and contraction of a surrounding particle sphere. Inhalation expands the environment and exhalation contracts it. Coupled oscillators govern relations among particle movements and appearance, allowing visual organization to be manipulated while retaining the same respiratory mapping. The controlled implementation runs on a standalone Meta Quest 3 headset. A controller positioned against the abdomen provides the movement signal; the scene applies smoothing to obtain stable visual feedback. The system papers document the development and a subsequent design extension integrating cardiac biofeedback. These engineering outputs do not by themselves establish physiological or clinical efficacy. [3,4,5]

\subsection{Exploratory pilot work}

The first progress report described an exploratory sample of 30 participants, comprising 20 visitors at Gamesground in Berlin and 10 University of Konstanz students. It documented feasibility, subjective experiential observations and retrospective temporal tracers. Those observations motivated subsequent controlled testing. They are retained here as pilot evidence rather than as a confirmatory test: the available progress report does not provide a complete inferential analysis supporting its stronger descriptions of condition differences. The pilot sample must not be combined with the later controlled sample as though they formed one experiment. [6]

\subsection{Participant accounting across project studies}

Table~\ref{tab:participant-accounting} consolidates participant numbers for each empirical strand. Recruitment targets, recruited or tested samples, exclusions, analysed samples and ongoing recruitment are reported separately. Counts from different studies are not pooled, because participant overlap has not been established and the studies address different outcomes.

{
\small\setlength{\tabcolsep}{4pt}\renewcommand{\arraystretch}{1.16}
\begin{longtable}{>{\raggedright\arraybackslash}p{0.24\linewidth}>{\raggedright\arraybackslash}p{0.28\linewidth}>{\raggedright\arraybackslash}p{0.40\linewidth}}
\caption{Participant accounting by empirical study, as documented on 30 October 2026.}\label{tab:participant-accounting}\\
\toprule
\textbf{Study / phase} & \textbf{Recruitment or records} & \textbf{Analysis status and interpretation}\\
\midrule\endfirsthead
\toprule
\textbf{Study / phase} & \textbf{Recruitment or records} & \textbf{Analysis status and interpretation}\\
\midrule\endhead
\bottomrule\endfoot
Exploratory feasibility pilot & 30 observed: 20 Gamesground Berlin visitors and 10 Konstanz students & Exploratory observations only; no verified confirmatory analysed sample or exclusion total in the available progress report.\\
Controlled subjective evaluation (Study 3 / Study II subjective strand) & Preregistered target 36 completers; 47 participants tested & 8 excluded by the reported phase-locking-value criterion; 39 participants included in the three-condition analyses.\\
Behavioural peripersonal-space study (Study 5 / Study II objective strand) & Public release: 42 participant records and 17,114 observed trial rows & Release-inventory counts, not a verified final eligible or analysed sample; participant QC and boundary-identifiability analyses remain ongoing.\\
Affective mapping and HRV study (Study 6) & Preregistered target 24; investigator-reported recruitment 8/24 & Data collection incomplete; no finalized exclusions or confirmatory analytical sample established.\\
Study I system development and PPS Kit & Technical and methodological deliverables; no separate participant cohort stated & Not additional empirical participant samples; do not add these outputs to participant totals.\\
\end{longtable}
}

The final submission should update the Study 5 and Study 6 rows if recruitment, quality-control decisions or analysed sample sizes change. No unique project-wide participant total is asserted without resolving overlap and eligibility.

\subsection{Study II subjective evaluation participants and procedure}

The controlled study tested 47 participants recruited from the University of Konstanz psychology participant pool. The preregistered recruitment target was 36 completers, balanced across six condition orders. Eight participants were excluded using the reported phase-locking-value criterion, leaving 39 participants with complete observations for all three conditions. The included sample had a mean age of 20.8 years and comprised 37 women and 2 men. Prior VR, meditation, breathwork and psychedelic experience was not systematically assessed. These figures describe the manuscript dataset, not the peripersonal-space study. [5]

Participants sat in a bean bag, wore headphones and held a controller against the abdomen. After a one-minute practice block, they completed three ten-minute audio-guided box-breathing sessions in one of six counterbalanced orders. Each cycle comprised four seconds of inhalation, breath hold, exhalation and breath hold. The conditions were a high-symmetry particle environment, an asymmetric particle environment, and a dark-field control displaying a minimal central breathing-performance indicator. The control therefore contained visual feedback and was not a no-feedback or eyes-closed condition. The two VR conditions differed in particle organization during breath holds, not in whether inhalation expanded or contracted space. [5]

\subsection{Instruments and analysis}

Outcomes comprised breathing adherence, three pictographic self-related measures, the 42-item 11D-ASC questionnaire and exploratory temporal experience tracers. The self-related measures assessed perceived body boundaries, spatial frame of reference and small self. Higher-order 3D-ASCr composites summarized positive, perceptual and distressing effects from the same ASC responses; they are not independent datasets. Breathing adherence was operationalized as phase-locking value between the recorded sphere-radius trajectory and an ideal trajectory derived from the audio instructions. [5]

The manuscript documents a change from the preregistered phase-duration measure to phase-locking value because controller-derived phase segmentation was unstable. A pooled lower-tail interquartile-range rule applied at the condition level excluded an entire participant if any condition was flagged. The reported planned comparisons were paired tests of pooled VR versus the control and symmetric versus asymmetric VR, with order-adjusted mixed models and nonparametric and false-discovery-rate sensitivity analyses. Statistics below are transcribed from the cited manuscript snapshot and require author confirmation against the final analysis outputs. The separate reproducibility handoff identifies a withdrawn post-unmasking endpoint; no claim based on that endpoint is included here. [5]

\authoraction{Confirm the ethics approval reference and dates for each implemented study. The manuscript states University of Konstanz ethics approval and informed consent; the precise amendment coverage should be checked against the archived approval documents.}

\subsection{Study 5: behavioural peripersonal-space investigation}

Study 5 investigates respiratory and intervention-related modulation of peripersonal-space representation. Its deidentified behavioural data, analysis code and evolving results are publicly available in the repository Respiratory Modulation of Peripersonal Space Representation. The released package contains 42 participant records and 17,114 observed trial rows; these are release-inventory counts rather than a verified final analytical sample. Data analysis remains ongoing because the response patterns have not yielded reproducible and reliably identifiable PPS boundaries under the assumed sigmoid-function model. [7]

\authoraction{Confirm final analytical cohorts and recording linkage, prespecified outcomes and deviations, signal/trial quality rules, and completed boundary-identifiability and sensitivity analyses. Keep the 42 release records separate from the 39-person subjective-study sample.}

\subsection{Study 6: affective mapping and physiological measures}

Study 6 separately examines whether affectively parameterised Viscereality scenarios elicit the intended changes in valence and arousal. It combines cardiac recordings and heart-rate-variability metrics with psychometric assessment of affective experience; HRV is treated as a physiological correlate to be interpreted alongside subjective measures. According to the investigator update of 30 October 2026, data collection has stalled at 8 of the planned 24 preregistered participants. Recruitment and data collection therefore remain incomplete, and no completed confirmatory result or validated affective-induction effect is claimed. The final report will need the remaining recruitment outcome, a reconciled account of the registered and implemented protocol, final signal-quality and exclusion decisions, and completed analyses before this strand can be regarded as concluded. [8]

\subsection{PPS Kit: replication framework}

The PPS replication tool developed within the project has evolved into PPS Kit, a functional framework for designing, preparing and running audio-tactile peripersonal-space paradigms. It consolidates stimulus-design parameters available from the reviewed literature into reusable profiles and a common design workflow, with source citations, field-level provenance and explicit missing-parameter flags. The documented library contains two Study 5 workflows and 28 published-study profile variants. PPS Designer provides the design interface and PPS Experiment Runner executes prepared experiments. This work builds on the measurement and reproducibility difficulties encountered in Study 5 and supports systematic replication and extension of the covered paradigms. It is a methodological output; successful software operation does not establish replication of every published effect, exact equivalence to every original apparatus, or a validated PPS boundary measure. [9]

Public toolkit and documentation: \url{https://ppskit.qzz.io/}.
